% Leitura geral do catálogo e da pesquisa. Texto editável (não gerado).

\subsection*{Resumo do catálogo}

\begin{itemize}
  \item \textbf{66 features em 12 módulos}: 10 no MVP, 21 na V2, 19 na V3, 9
    na V4, 5 na trilha condicional e 2 fora (tabela~\ref{tab:fases}). São 59
    features comprometidas, 5 em opção e 2 descartadas. 47 têm encaixe alto com o mapa, 12
    médio e 7 baixo.
  \item \textbf{A distribuidora concentra o valor espacial}: 50 das 66
    features servem a ela (tabela~\ref{tab:clientes}). Geração renovável e
    mercado têm encaixe médio ou baixo, o que é coerente com a hipótese de cliente
    da especificação.
  \item \textbf{CON-01 e CON-02 estão na V2}: a BDGD traz DIC e FIC mensais
    por unidade. CON-02 (risco de violação de DIC e FIC) é de nível 0, tem
    encaixe alto e mede a compensação, em reais, paga pela distribuidora.
    CON-01 é pré-requisito de CON-02.
  \item \textbf{O MVP fecha o ciclo}: base (BAS-01 a 03), perdas
    (PER-01 a 04) e engine (ENG-01 a 03) cobrem feed $\rightarrow$
    diagnóstico $\rightarrow$ plano $\rightarrow$ resultado.
\end{itemize}

\subsection*{Efeitos da pesquisa sobre o catálogo}

\paragraph{1. Dados disponíveis no nível 0.} Com dado público
apenas, existem:
\begin{itemize}
  \item na BDGD: \textbf{energia mensal por alimentador, perdas técnicas e não
    técnicas por alimentador}, energia, DIC e FIC mensais por unidade, e o
    modelo elétrico que a ANEEL roda no OpenDSS;
  \item na ANEEL: \textbf{cada interrupção desde 2017 com o código do
    alimentador e da subestação}, causa e expurgo; DEC e FEC com limites;
  \item no ONS: geração por usina, constrained-off semi-horário por usina com
    razão e ponto de conexão, carga com a parcela de MMGD.
\end{itemize}
Por isso, 36 das 66 features têm versão em nível 0, entre elas PER-01 (por
alimentador), CAR-02, RED-01, CON-01 a 03, GER-01 e GER-04; quatro delas
(MMG-06, RED-01, RED-02 e ATV-02) dependem de uma feature de nível 1 ou 2. \textbf{Um
piloto em modo sombra pode ser montado antes de receber qualquer arquivo do
cliente.}

\paragraph{2. Limites do nível 0.}
\begin{itemize}
  \item \textbf{Defasagem de 8 a 20 meses} na BDGD (edição de 2025): serve para montar a rede e
    demonstrar, não para operar mês a mês.
  \item \textbf{A unidade consumidora é anonimizada} (hash) na BDGD pública.
    Ligar o dado do cliente por unidade exige a tabela de correspondência da
    distribuidora; pela rede (alimentador, transformador) a ligação tende a
    funcionar.
  \item \textbf{Unidades inconsistentes entre distribuidoras} (kWh e MWh no
    mesmo campo) e, na MMGD, injeção no lugar de geração e kW misturado com W,
    observados em duas BDGD de 2025. A normalização por distribuidora e a
    camada de qualidade (BAS-03) são necessárias.
  \item A perda não técnica das duas BDGD é compatível com \textbf{rateio
    contábil}, não com medição.
\end{itemize}

\paragraph{3. Três regras de validação atravessam o catálogo.}
\begin{itemize}
  \item \textbf{Guardar o que era conhecido em cada data}: ONS e CCEE revisam
    dados depois de publicados, e a previsão do tempo do ECMWF só fica
    disponível por 2 a 3 dias. Arquivar desde o primeiro dia é requisito.
  \item \textbf{Validar fora do espaço e do evento}: em interrupções por clima,
    divisões aleatórias inflam a estimativa de desempenho; em Essus et
    al.\ (2026), com validação espacial e temporal, muitos modelos não superam
    a linha de base.
  \item \textbf{Controlar testes múltiplos} na engine: entidades × camadas ×
    detectores exigem FDR e orçamento de alertas.
\end{itemize}

\paragraph{4. Fundamento da fração exploratória.} Revelas et al.\ (2025)
mostram que selecionar só os alvos mais prováveis leva a aprendizado
inconsistente; a seleção aleatorizada com probabilidades registradas permite
estimar, sem viés, prevalência, recall, calibração e o viés de seleção. A mesma fração alimenta PER-02, PER-05 e PER-06.

\paragraph{5. Regras regulatórias estão mudando.}
\begin{itemize}
  \item Perda não técnica calculada sobre o mercado medido desde 2025, o que
    quebra a série.
  \item Formação de preço por oferta e dois PLDs previstos para 2028 (Meta II).
  \item REN 1.098/2024 sobre inversão de fluxo da MMGD.
  \item Portaria MME 126/2026: 2\% ao ano de medidores inteligentes,
    priorizando perdas.
  \item Proposta de Tarifa Branca automática.
\end{itemize}
Cada uma é um ponto de deriva a monitorar (ENG-04).

\paragraph{6. Licenças a observar.} BDGD em ODbL (compartilhamento pela mesma
licença em bases derivadas públicas: consultar o jurídico); ECMWF em CC-BY com
atribuição no mapa; \emph{alibi-detect} em Business Source License, com
restrição comercial.

\subsection*{A trilha condicional}

Cinco features --- MER-01 a MER-03, CAR-08 e GER-08 --- saíram do roadmap
faseado para a trilha \textbf{cond}. Não são features descartadas: são uma
opção sobre a questão de cliente-alvo que a especificação deixa em aberto
(distribuidora, ou gerador e comercializadora). Enquanto essa questão não for
respondida, elas não recebem fase.

A retirada não afeta as demais: formam uma subárvore isolada
---~CAR-08 e GER-08 $\rightarrow$ MER-01 $\rightarrow$ MER-02 $\rightarrow$
MER-03~--- e \textbf{nada fora dela depende de nenhuma das cinco}. As
dependências externas da subárvore são CLI-01 (de CAR-08) e GER-01 (de GER-08
e MER-02), features que permanecem no roadmap.

Três observações pesaram na decisão:
\begin{itemize}
  \item \textbf{É a parte do gerador em que o mapa não ajuda.} As cinco têm
    encaixe baixo, e o preço por submercado ocupa quatro polígonos. Se a trilha
    de gerador for tomada, a âncora espacial é \textbf{RED-04} (restrição e
    curtailment por região): encaixe alto, nível 0, sem dependências, e já
    sustentada pelo constrained-off semi-horário do ONS.
  \item \textbf{A premissa de MER-03 só existe a partir de 2028.} Oferta só faz sentido com
    o modelo híbrido do Projeto Meta II, e o primeiro leilão de baterias
    (Portaria MME 136/2026) remunera disponibilidade, não arbitragem livre.
  \item \textbf{A continuidade da fonte não é garantida.} O portal de dados
    abertos da CCEE tem bloqueio de aplicação; o endpoint CKAN funcionou em
    23/09/2026.
\end{itemize}

\subsection*{Cadeias de dependência}

As dependências do catálogo definem a ordem de construção. As principais do roadmap faseado:
\begin{itemize}
  \item BAS-01 $\rightarrow$ PER-01, PER-02 $\rightarrow$ PER-03 $\rightarrow$
    PER-04 $\rightarrow$ PER-05, PER-06 (o MVP e sua avaliação);
  \item MMG-01 $\rightarrow$ MMG-02 $\rightarrow$ MMG-03, MMG-06, MMG-07;
  \item CAR-01, CAR-02 $\rightarrow$ RED-01 $\rightarrow$ RED-03, DRE-01;
  \item ATV-03 $\rightarrow$ CON-03, e ATV-01, ATV-02 $\rightarrow$ ATV-06.
\end{itemize}

\subsection*{Pendências de verificação}

\begin{enumerate}
  \item Manual de Instruções da BDGD (não baixado): definição oficial de
    energia e perdas na CTMT, energia das unidades geradoras (injeção ou
    geração), unidade da potência instalada.
  \item Se o código de conjunto da BDGD casa com o dos indicadores de DEC e
    FEC.
  \item Se a distribuidora consegue reproduzir o hash do código da unidade.
  \item Unidade do valor nas curvas típicas CTR.
  \item Existência de obrigação regulatória de publicar capacidade de
    hospedagem.
  \item Consulta jurídica sobre a ODbL em camadas derivadas.
\end{enumerate}
